\documentclass{zkdl-presentation-template}

\title[NTT]{\textbf{Number Theoretic Transform (NTT)}}
\author{Blockstream Research}
\date{\today}
\homepage{research.blockstream.com}
\github{BlockstreamResearch/zk-book}

\begin{document}
    \maketitleframe

    \begin{frame}{Plan}
        \tableofcontents
    \end{frame}

    \section{Recap on Interpolation}

    \subsection{Polynomial Interpolation is a Universal Encoder}
    \begin{frame}{Polynomial Interpolation}
        \begin{remarkblock}{Notice}
            All the previous protocols use the idea that polynomials
            are \textbf{universal data encoders}.\pause{} We can encode any
            set of scalars $(a_0,\dots,a_{N-1}) \in \mathbb{F}^N$
            using \textbf{interpolation}:
            \begin{equation*}
                p(x_j) = a_j, \quad j \in [N], \quad \text{$\{x_j\}_{j \in [N]}$ are fixed}
            \end{equation*}
        \end{remarkblock}

        \begin{figure}
            \centering
            \includegraphics[width=\textwidth]{images/lecture_13/encoding.pdf}
            \caption{Polynomial Interpolation as a universal encoder.}
        \end{figure}
    \end{frame}

    \begin{frame}{Polynomial Interpolation}
        \begin{example}
            In Groth16, we used interpolation of $3n$ polynomials:
            \begin{equation*}
                \ell_j(i) = L_{i,j}, \quad r_j(i) = R_{i,j}, \quad o_j(i) = O_{i,j},
            \end{equation*}
            where $L_{i,j}, R_{i,j}, O_{i,j}$ are the elements of constraint
            matrices $L, R, O$ (left, right, and output).
        \end{example}
        \pause

        \begin{columns}
            \begin{column}{0.65\textwidth}
                However, in PlonK (see the book) we will interpolate $a(\omega^j) = A_j$ where
                $A_j$ are the elements of the left trace vector $A$.

                \begin{questionblock}{Question}
                    What the heck is this $\omega$? Why do we
                    need it? How does it help?
                \end{questionblock}
            \end{column}
            \begin{column}{0.35\textwidth}
                \begin{figure}
                    \centering
                    \includegraphics[width=0.8\textwidth]{images/lecture_13/plonk_table.pdf}
                \end{figure}
            \end{column}
        \end{columns}
    \end{frame}

    \subsection{Motivation for NTT}
    \begin{frame}{Why we need something advanced?}
        \begin{remarkblock}{Recall}
            The interpolation formula is given by:
            \begin{equation*}
                p(T) = \sum_{i=0}^{N-1} a_i \cdot \ell_i(T), \quad \ell_i(T) = \prod_{j=0, j \neq i}^{N-1} \frac{T - x_j}{x_i - x_j}
            \end{equation*}
        \end{remarkblock}
        \pause

        \begin{questionblock}{Question}
            What is the naive complexity of this interpolation implementation?
        \end{questionblock}
        \pause

        \begin{block}{Observation}
            Through careful choice of $\{x_j\}_{j \in [N]}$, we can reduce the
            complexity of interpolation, multiplication, or other complex
            operations to $O(N \log N)$.\pause{} \textbf{Spoiler:} we will use
            the $N$th roots of unity domain $\Omega = \{\omega^j\}_{j \in [N]}$.
            Let us see why it helps.
        \end{block}
    \end{frame}

    \section{Roots of Unity}
    \subsection{Multiplicative Subgroup of Finite Fields}
    \begin{frame}{Multiplicative Subgroup}
        We know that $\mathbb{F}_p$ is a \textbf{field}: we have a usual arithmetic $+,\times$.
        \pause

        \begin{questionblock}{Question}
            Does $(\mathbb{F}_p, \times)$ form a group?
        \end{questionblock}
        \pause

        No, since $0$ does not have an inverse. But, if we consider
        $(\mathbb{F}_p \setminus \{0\}, \times)$, we do have a group structure!
        \pause

        \begin{definition}
            A \textbf{multiplicative group} of a finite field $\mathbb{F}$, denoted as $\mathbb{F}^{\times}$, is a multiplicative group $(\mathbb{F} \setminus \{0\}, \times)$.
        \end{definition}
        \pause

        \begin{remarkblock}{Number of Elements}
            The number of elements in $\mathbb{F}^{\times}_p$ is $p - 1$.
        \end{remarkblock}
    \end{frame}

    \begin{frame}{Primitive Root}
        \begin{theorem}
            Multiplicative group of a finite field $\mathbb{F}^{\times}_p$ is
            \textbf{cyclic}. The generators $\omega$ of this group are called
            \textbf{primitive roots}.
        \end{theorem}
        \pause

        \begin{example}
            $\omega=3$ is the primitive root of $\mathbb{F}_7$. Indeed,
            \begin{align*}
                3^1 = 3, \quad 3^2 = 2, \quad 3^3 = 6, \quad 3^4 = 4, \quad 3^5 = 5, \quad 3^6 = 1.
            \end{align*}

            Clearly, $\langle \omega \rangle = \mathbb{F}_7^{\times}$.
        \end{example}
        \pause

        The set $\mathbb{F}^{\times}_p$ is not useful on its own. However, we
        can consider the following set, called \textbf{$r$-th roots of unity}:
        \begin{equation*}
            \Omega_r = \{ \omega \in \mathbb{F}^{\times}_p : \omega^r = 1 \} \subset \mathbb{F}_p^{\times}.
        \end{equation*}

        \vspace{-15px}
        \pause

        \textcolor{oc-pink-7}{\textbf{Question.}} When does such a cyclic group exist?
    \end{frame}

    \begin{frame}{Roots of Unity}
        \begin{theorem}[Lagrange Theorem]
            If $\mathbb{H} \leq \mathbb{G}$ is a subgroup of any finite group $\mathbb{G}$, then $\#\mathbb{H} \mid \#\mathbb{G}$.
        \end{theorem}
        \pause

        \begin{corollary}
            If $\Omega_r$ is a subgroup of $\mathbb{F}_p^{\times}$, then $r \mid
            (p - 1)$.
        \end{corollary}
        \pause

        \begin{remarkblock}{Some other Notes}
            Moreover, one might prove in the opposite direction:
            \begin{itemize}
                \item If $r \mid (p-1)$, then there exists a subgroup $\Omega_r
                \leq \mathbb{F}_p^{\times}$.\pause
                \item Its generator is given by $\omega = g^{(p-1)/r}$ where
                $\langle g \rangle = \mathbb{F}_p^{\times}$.
            \end{itemize}
        \end{remarkblock}
        \pause

        \begin{remarkblock}{Yet another note}
            Typically, we would need $r$ to be the power of two. We will see why
            in the NTT section.
        \end{remarkblock}
    \end{frame}

    \begin{frame}{Complex Analysis Interpretation}
        \begin{figure}
            \centering
            \includegraphics[width=0.5\textwidth]{images/lecture_13/roots.png}
            \caption{Visualization of the roots of unity $\Omega_5 = \{z \in \mathbb{C}: z^5=1 \}$.}
        \end{figure}

        On the complex plane, the generator of the $r$-th roots of unity
        $\Omega_r$ is given by $\zeta_r = e^{2\pi i/r}$. In a finite field,
        we do not have such a luxury.
    \end{frame}

    \subsection{Barycentric Interpolation}
    \begin{frame}{Vanishing Polynomial}
        \begin{definition}
            The \textbf{vanishing polynomial} $z_D(T)$ of a set $D \subset \mathbb{F}_p$
            is a polynomial satisfying $z_D(d) = 0$ for all $d \in D$.
        \end{definition}
        \pause
        Vanishing polynomials are always of form $z_D(T) = c \cdot \prod_{d \in
        D}(T-d)$.
        \pause

        The interesting question is: what is the vanishing polynomial of the
        $r$-th roots of unity $\Omega_r$? For simplicity, assume $c=1$.
        \pause

        \begin{lemma}
            The vanishing polynomial of $\Omega_r$ is $z_{\Omega}(T) = T^r-1$.
        \end{lemma}
        \pause

        \textbf{Proof Idea.} Since for any $\zeta \in \Omega_r$ we have $\zeta^r=1$, or,
        equivalently, $\zeta^r-1=0$.\pause{} Thus, any $\zeta \in \Omega_r$ is a root of $z_{\Omega}(T) = T^r-1$.
    \end{frame}

    \begin{frame}{Vanishing Polynomial over $\mathbb{R}$}
        \begin{figure}
            \centering
            \includegraphics[width=0.7\textwidth]{images/lecture_13/vanishing.png}
            \caption{Vanishing polynomial $p(x)=(x-1)(x-2)(x-4)$ of $D=\{1,2,4\}$}
        \end{figure}
    \end{frame}

    \begin{frame}{Barycentric Interpolation}
        Now, let us come back to the interpolation problem $p(x_j) = a_j$ for $j \in [N]$.
        Introduce $\gamma(T) = \prod_{j=0}^{N-1}(T-x_j)$.
        \pause

        \begin{proposition}
            The Lagrange basis polynomial $\ell_j$ can be rewritten as:
            \begin{equation*}
                \ell_j(T) = \gamma(T) \cdot \frac{w_j}{T-x_j}, \quad w_j = \frac{1}{\prod_{k=0,k\neq j}^{N-1}(x_j-x_k)}.
            \end{equation*}
        \end{proposition}
        \pause

        Let us substitute it into the interpolation formula:
        \begin{equation*}
            p(T) = \sum_{j=0}^{N-1}a_j\ell_j(T) = \sum_{j=0}^{N-1}a_j\gamma(T)\frac{w_j}{T-x_j} \pause = \gamma(T)\sum_{j=0}^{N-1}\frac{w_j}{T-x_j}a_j.
        \end{equation*}
    \end{frame}

    \begin{frame}{Barycentric Interpolation (Cont.)}
        \begin{equation*}
            \boxed{\textbf{Barycentric Formula:} \; \; p(T) = \gamma(T)\sum_{j=0}^{N-1}\frac{w_j}{T-x_j}a_j}
        \end{equation*}
        \pause

        \begin{proposition}
            \begin{itemize}
                \item Computing $\{w_j\}_{j \in [N]}$ costs $O(N^2)$
                operations \emph{before evaluation}.\pause
                \item Both $\gamma(T)$ and sum require $O(N)$ operations.
            \end{itemize}
        \end{proposition}
        \pause

        But what happens if instead of $x_j$, we use $\omega^j \in \Omega_N$?\pause
        \begin{equation*}
            \boxed{p(T) = \frac{T^N-1}{N}\sum_{j \in [N]}\frac{\omega^j}{T-\omega^j}a_j}
        \end{equation*}
        \pause

        \textcolor{oc-pink-7}{\textbf{Takeaway:}} We can interpolate in $O(N)$ operations.
    \end{frame}

    \section{Number Theoretic Transform}

    \subsection{Three Gadgets}
    \begin{frame}{What is NTT?}
        Now suppose we want to find $m(T) = p(T)q(T)$. We'll use NTT!
        \pause

        \begin{questionblock}{Question}
            What does it mean that you \emph{know} polynomial $p(T) \in \mathbb{F}[T]^{< N}$?
        \end{questionblock}
        \pause

        This means either of two (typically):
        \begin{itemize}
            \item You know the polynomial coefficients $p_0,\dots,p_{N-1}$.\pause
            \item You know the polynomial values at some points $\{(x_j,a_j)\}_{j \in [N]}$.
        \end{itemize}
        \pause

        \begin{definition}[NTT]
            Suppose $p(T) = \sum_{j=0}^{N-1}p_jT^j$. The \textbf{Number Theoretic Transform} (NTT) of $p$ is defined as evaluations of $p$ at the $N$-th roots of unity:
            \begin{equation*}
                \mathsf{NTT}(\mathbf{p}) = \left(p(\omega^0), p(\omega^1), \dots, p(\omega^{N-1})\right).
            \end{equation*}
        \end{definition}
    \end{frame}

    \begin{frame}{What is the point of NTT?}
        \textcolor{oc-green-8}{\textbf{Note:}} To denote the result of NTT, we use hat: $\hat{\mathbf{p}} = \mathsf{NTT}(\mathbf{p})$.
        \pause

        \textcolor{oc-pink-7}{\textbf{Question:}} Given NTTs $\hat{\mathbf{p}}$ and $\hat{\mathbf{q}}$ of two polynomials $p$ and $q$,
        how do we find the NTT of their product $m(T) = p(T)q(T)$?
        \pause

        \begin{proposition}[Main NTT Property]
            Suppose $m(T) = p(T)q(T)$ is the product of $p$ and $q$. Then,
            \begin{equation*}
                \hat{\mathbf{m}} = \hat{\mathbf{p}} \odot \hat{\mathbf{q}}
            \end{equation*}
            \pause

            Speaking more formally, $\mathsf{NTT}: (\mathbb{F}[T]^{\leq N},\times) \to (\mathbb{F}^N, \odot)$ is
            a homomorphism between a set of polynomials of degree up to $N$ and their NTT domain. With certain
            appropriate technicalities, $\mathsf{NTT}$ can be extended to the isomorphism (namely, use $\mathbb{F}[T] / (T^N-1)$).
        \end{proposition}
        \pause

        \textcolor{oc-blue-7}{\textbf{Why?}}\pause{} Well... $m(\omega^j)=p(\omega^j)q(\omega^j)$ \hspace{20px} \textbf{:/}
    \end{frame}

    \begin{frame}{Final Ingredient: Inverse NTT}
        Now, can we restore the polynomial $m(T)$ from its NTT $\hat{\mathbf{m}}$?\pause{} Of course!

        \begin{definition}[Inverse NTT]
            The \textbf{Inverse Number Theoretic Transform} (INTT) is a function
            that restores the polynomial $m(T)$ from its evaluations $\hat{\mathbf{m}}$:
            \begin{equation*}
                \mathsf{INTT}(\hat{\mathbf{m}}) = (m_0,m_1,\dots,m_{N-1})
            \end{equation*}
        \end{definition}
        \pause

        In its essence, we solve the interpolation problem:
        \begin{equation*}
            m(\omega^j) = \hat{m}_j, \quad j \in [N], \quad \text{\textbf{Goal}: find coefficients $m_0,\dots,m_{N-1}$}
        \end{equation*}
    \end{frame}

    \subsection{Polynomial vs NTT Domain}

    \begin{frame}{Punchline}
        \begin{figure}[H]
            \centering
            \begin{tikzpicture}[node distance=2.5cm, >=Stealth]
                % Polynomial zone
                \node[very thick, draw=oc-blue-7, fill=oc-blue-1, dashed, minimum width=9.5cm, minimum height=1.75cm, on background layer](polybox) at (2.25, 0.1) {};
                \node[above=0.65cm of polybox, anchor=north, color=oc-blue-8]{\textbf{Polynomial Space $\mathbb{F}[T]^{\leq N}$}};

                \node (pq) at (0,0) {\( p, q \)};
                \node (m) [right=3cm of pq] {\( m = p \cdot q \)};
                \draw[->, line width = 0.5mm] (pq) -- (m) node[midway, above] {$O(N^2)$};
                \pause

                % NTT zone
                \node[very thick, draw=oc-green-7, fill=oc-green-1, dashed, minimum width=9.5cm, minimum height=2.0cm, on background layer](nttbox) at (2.25, -2.0) {};
                \node[below=0.15cm of nttbox, anchor=north, color=oc-green-8]{\textbf{NTT Space}};
                \pause

                \node (pHatqHat) [below=2cm of pq] {\( \hat{\mathbf{p}}, \hat{\mathbf{q}} \)};
                \node (mHat) [below=2cm of m] {\( \hat{\mathbf{m}} = \hat{\mathbf{p}} \odot \hat{\mathbf{q}} \)};
                \pause

                \draw[->, dashed, line width = 0.5mm] (pq) -- (pHatqHat) node[midway, left] {$O(N \log N)$};
                \draw[->, dashed, line width = 0.5mm] (pHatqHat) -- (mHat) node[midway, above] {$O(N)$};
                \draw[->, dashed, line width = 0.5mm] (mHat) -- (m) node[midway, right] {$O(N \log N)$};
            \end{tikzpicture}
            \caption{Illustration of the NTT Algorithm}
        \end{figure}
        \pause

        \begin{questionblock}{Question}
            Does it remind you of one trick from Elliptic Curves?
        \end{questionblock}
    \end{frame}

    \begin{frame}{Illustration}
        \begin{figure}
            \centering
            \includegraphics[width=\textwidth]{images/lecture_13/ntt_visualization.png}
            \caption{Illustration of the FFT Algorithm. Taken from ``The Fast Fourier Transform (FFT): Most Ingenious Algorithm Ever?''}
        \end{figure}
    \end{frame}

    \section{Details}

    \begin{frame}{When NTT works?}
        \begin{remarkblock}{Note}
            For NTT to work, we will impose the following requirements on
            our setup:
            \begin{enumerate}
                \item The field $\mathbb{F}_p$ should have $2^k$-roots of unity for sufficiently many $k$. In other words, $p=p' \cdot 2^m + 1$ with \emph{small} $p' \in \mathbb{N}$.\pause
                \item The polynomial order is $N=2^k$. Not a strict requirement, since we can always pad the polynomial with zeros.
            \end{enumerate}
        \end{remarkblock}
        \pause

        \begin{example}
            \begin{itemize}
                \item \textbf{BabyBear prime} $p=15 \cdot 2^{27} + 1$ is NTT-friendly: the order of multiplicative
                group is $15 \cdot 2^{27}$, so $2^k \mid 15 \cdot 2^{27}$ for all $k \leq 27$.\pause
                \item \textbf{Mersenne prime} $p=2^{31}-1$ is not NTT-friendly: the order of multiplicative
                group is $2^{31}-2 = 2 \times (2^{30}-1)$.
            \end{itemize}
        \end{example}
    \end{frame}

    \subsection{Why NTT takes quasilinear complexity?}

    \begin{frame}{Why NTT takes quasilinear complexity?}
        Recall that we need to evaluate $N$ expressions:
        \begin{equation*}
            p(\omega^j) = \sum_{i=0}^{N-1}p_i(\omega^j)^i = \sum_{i=0}^{N-1}p_i\omega^{ij}, \quad j \in [N].
        \end{equation*}
        \pause

        \textcolor{oc-pink-7}{\textbf{Naive Complexity:}} $O(N^2)$ operations. We need $N$ evaluations,
        each of which requires $N$ multiplications.
        \pause

        \begin{align*}
            p(\omega^j) &= \sum_{i=0}^{2^r-1}p_i\omega^{ij} = \sum_{i=0}^{2^{r-1}-1}p_{2i}\omega^{2ij} + \sum_{i=0}^{2^{r-1}-1}p_{2i+1}\omega^{j(2i+1)}\\
            &= \textcolor{oc-green-8}{\sum_{i=0}^{2^{r-1}-1}p_{2i}(\omega^{2j})^i} + \omega^j\textcolor{oc-blue-7}{\sum_{i=0}^{2^{r-1}-1}p_{2i+1}(\omega^{2j})^i}.
        \end{align*}
    \end{frame}

    \begin{frame}{Folding}
        Denote $p_E(T) = \sum_{i=0}^{2^{r-1}-1}p_{2i}T^i$ and $p_O(T) = \sum_{i=0}^{2^{r-1}-1}p_{2i+1}T^i$. Then,
        \begin{equation*}
            p(\omega^j) = \textcolor{oc-green-8}{p_E(\omega^{2j})} + \omega^j \textcolor{oc-blue-7}{p_O(\omega^{2j})}.
        \end{equation*}
        \pause

        \begin{block}{Fact \#1}
            We need only $N/2$ evaluations from $\Omega$ of $p_E$ and $p_O$. Note that:
            \begin{equation*}
                p(\omega^{j+N/2}) = \textcolor{oc-green-8}{p_E(\omega^{2j})} + \omega^j\omega^{N/2} \textcolor{oc-blue-7}{p_O(\omega^{2j})}.
            \end{equation*}
        \end{block}
        \pause

        \begin{block}{Fact \#2}
            \begin{itemize}
                \item We need to evaluate two $N/2$-degree polynomials.\pause
                \item We need to evaluate them at $N/2$ points.
            \end{itemize}
            \pause

            Thus, we shrink the problem size by half at each step.
        \end{block}
    \end{frame}

    \begin{frame}{Algorithm Summarized}
        \begin{algorithm}[H]
            \caption{Number Theoretic Transform (NTT)}
            \Input{Polynomial $p(T) = \sum_{j=0}^{N-1}p_jT^j$}
            \Output{Vector of evaluations $\mathsf{NTT}(\mathbf{p}, \omega)$ at $\Omega=\{\omega^j\}_{j \in [N]}$}

            \If{$N=1$}{
                \Return{$(p_0)$}
            }

            $H \gets N/2$ \Comment{Compute the domain half-size}

            $\mathbf{p}_{E} \gets (p_0,p_2,\dots,p_{N-2})$ \Comment{Find even-indexed coefficients}

            $\mathbf{p}_{O} \gets (p_1,p_3,\dots,p_{N-1})$ \Comment{Find odd-indexed coefficients}

            $\mathbf{y}_E \gets \mathsf{NTT}(\mathbf{p}_{E}, \omega^2)$ \Comment{Compute NTT for even polynomial via $\frac{N}{2}$th primitive root $\omega^2$}

            $\mathbf{y}_O \gets \mathsf{NTT}(\mathbf{p}_O, \omega^2)$ \Comment{Compute NTT for odd polynomial via $\frac{N}{2}$th primitive root $\omega^2$}

            \Return{$(y_0,\dots,y_{N-1})$ with $y_j=y_{E, \;j \; \text{mod} \; H} + \omega^j y_{O, \; j \; \text{mod} \; H}$}
        \end{algorithm}
    \end{frame}

    \begin{frame}{Inverse NTT}
        \begin{theorem}
            The Inverse NTT can be computed in the same way as NTT, but with the
            inverse primitive root $\omega^{-1}$:
            \begin{equation*}
                p_j = \frac{1}{N}\sum_{i=0}^{N-1}\omega^{-ij}\hat{p}_i
            \end{equation*}

            Thus, its complexity is also $O(N \log N)$.
        \end{theorem}
        \pause

        \begin{remarkblock}{Conclusion}
            To compute $m(T) = p(T)q(T)$, simply use the following:
            \begin{equation*}
                \mathbf{m} = \mathsf{INTT}(\mathsf{NTT}(\mathbf{p}) \odot \mathsf{NTT}(\mathbf{q}))
            \end{equation*}

            The total complexity remains $O(N \log N)$.
        \end{remarkblock}
    \end{frame}

    \makeclosingframe
\end{document}
